\documentclass[14pt,a4paper,oneside]{extbook}
\usepackage[utf8]{vietnam}
\usepackage{amsmath,amssymb,mathrsfs,fancyhdr,enumerate,multirow,makecell}
\usepackage{tkz-euclide,tikz-3dplot,tikz,tkz-tab}
\usetikzlibrary{arrows.meta,calc,intersections,angles,quotes,backgrounds,shapes.geometric,patterns,through,shadows,positioning}
\usepackage{yhmath}
\usepackage[top=0.7cm, bottom=0.7cm, left=1.0cm, right=1.0cm, headheight=16pt, headsep=3mm, footskip=5mm]{geometry}
\usepackage[hidelinks,unicode]{hyperref}
\usepackage{qrcode}
\usepackage[many]{tcolorbox}
\usepackage{fontawesome}

\pagestyle{fancy}
\fancyhf{}
\fancyhead[L]{\footnotesize\bfseries\color{teal!85!black}\faBook~SỔ TAY KIẾN THỨC TOÁN 5}
\fancyhead[R]{\footnotesize\bfseries\color{blue!85!black}Lớp Toán Thầy Lĩnh -- loptoanthaylinh.vn}
\fancyfoot[C]{\footnotesize\bfseries\color{gray!80!black}-- Sổ tay kiến thức: Bài 36. Tỉ số. Tỉ số phần trăm --}
\renewcommand{\headrulewidth}{0.4pt}

\begin{document}

% --- BANNER TIÊU ĐỀ BÀI HỌC ---
\noindent
\begin{tcolorbox}[
  enhanced,
  colback=teal!8!white,
  colframe=teal!80!black,
  arc=3mm,
  boxrule=1.1pt,
  width=\linewidth,
  halign=center,
  top=2.2mm, bottom=2.2mm,
  left=3mm, right=3mm,
  drop fuzzy shadow
]
  {\large\bfseries\color{teal!95!black}\faBookmark\ BÀI 36. TỈ SỐ. TỈ SỐ PHẦN TRĂM}\\[2pt]
  {\scriptsize\color{black!75}\itshape Khái niệm tỉ số a:b, ý nghĩa tỉ số phần trăm và cách chuyển đổi qua lại trên 1 trang A4}
\end{tcolorbox}

\vspace{3.2mm}
% --- KHỐI 1: KHÁI NIỆM TỈ SỐ & TỈ SỐ PHẦN TRĂM ---
\noindent
\begin{minipage}[t]{0.50\textwidth}
  \begin{tcolorbox}[
    enhanced,
    title={\footnotesize\bfseries\color{white}\faSliders\ 1. Tỉ số của hai số},
    colback=blue!2!white,
    colframe=blue!75!black,
    arc=2.5mm, boxrule=0.8pt, top=1.8mm, bottom=1.8mm, left=2.2mm, right=2.2mm
  ]
    \footnotesize
    $\bullet$ Tỉ số của số $a$ và số $b$ ($b \neq 0$) là thương $a : b$ hay $\dfrac{a}{b}$.\\
    $\bullet$ \textbf{Lưu ý thứ tự}: Tỉ số của $a$ và $b$ thì $a$ viết trước (ở tử số), $b$ viết sau (ở mẫu số).\\
    \textit{VD:} Có $3$ bạn nam và $5$ bạn nữ:\\
    $\triangleright$ Tỉ số giữa nam và nữ là: $3 : 5$ hay $\dfrac{3}{5}$.\\
    $\triangleright$ Tỉ số giữa nữ và cả lớp là: $5 : 8$ hay $\dfrac{5}{8}$.
  \end{tcolorbox}
\end{minipage}%
\hfill
\begin{minipage}[t]{0.48\textwidth}
  \begin{tcolorbox}[
    enhanced,
    title={\footnotesize\bfseries\color{white}\faPercent\ 2. Tỉ số phần trăm},
    colback=teal!2!white,
    colframe=teal!75!black,
    arc=2.5mm, boxrule=0.8pt, top=1.8mm, bottom=1.8mm, left=2mm, right=2mm
  ]
    \footnotesize
    $\bullet$ Là tỉ số có mẫu số bằng \textbf{100}.\\
    $\bullet$ Kí hiệu phần trăm là: \textbf{\%}.\\
    $\bullet$ $\dfrac{a}{100} = a\%$.\\
    \textit{VD:}\\
    $\triangleright$ $\dfrac{25}{100} = 25\%$ (hai mươi lăm phần trăm).\\
    $\triangleright$ $\dfrac{1}{2} = \dfrac{50}{100} = 50\%$.\\
    $\triangleright$ $\dfrac{3}{4} = \dfrac{75}{100} = 75\%$; $\dfrac{1}{4} = 25\%$.
  \end{tcolorbox}
\end{minipage}

\vspace{3mm}

% --- KHỐI 2: Ý NGHĨA THỰC TIỄN CỦA TỈ SỐ PHẦN TRĂM ---
\noindent
\begin{tcolorbox}[
  enhanced,
  title={\footnotesize\bfseries\color{white}\faEye\ 3. Ý nghĩa thực tiễn của Tỉ số phần trăm},
  colback=purple!2!white,
  colframe=purple!70!black,
  arc=2.5mm, boxrule=0.8pt, top=1.8mm, bottom=1.8mm, left=2.5mm, right=2.5mm
]
  \footnotesize
  Tỉ số phần trăm giúp so sánh các đại lượng trên cùng một thang đo chuẩn là \textbf{100 phần}:\\
  $\bullet$ \textit{Năng suất lúa tăng 15\%}: Nghĩa là cứ $100\text{ kg}$ vụ trước thì vụ này tăng thêm $15\text{ kg}$ (đạt $115\text{ kg}$).\\
  $\bullet$ \textit{Dung dịch muối 5\%}: Nghĩa là trong $100\text{ g}$ nước muối có đúng $5\text{ g}$ muối tinh khiết.\\
  $\bullet$ \textit{Cửa hàng giảm giá 20\%}: Nghĩa là món đồ $100\text{ nghìn}$ được giảm $20\text{ nghìn}$ (chỉ trả $80\text{ nghìn}$).
\end{tcolorbox}

\vspace{3mm}

% --- KHỐI 3: SƠ ĐỒ TƯ DUY ---
\noindent
\begin{tcolorbox}[
  enhanced,
  title={\footnotesize\bfseries\color{white}\faSitemap\ Sơ đồ tư duy: Chuyển đổi giữa Phân số -- STP -- Tỉ số phần trăm},
  colback=yellow!2!white,
  colframe=teal!85!black,
  arc=3mm, boxrule=0.9pt, top=1.8mm, bottom=1.8mm, left=1.5mm, right=1.5mm,
  halign=center
]
  \begin{tikzpicture}[
    scale=0.92,
    every node/.style={transform shape},
    b1/.style={rectangle, rounded corners=4pt, fill=blue!15, draw=blue!70!black, text=blue!90!black, font=\bfseries\scriptsize, inner sep=4pt},
    b2/.style={rectangle, rounded corners=4pt, fill=teal!15, draw=teal!70!black, text=teal!90!black, font=\bfseries\scriptsize, inner sep=4pt},
    b3/.style={rectangle, rounded corners=4pt, fill=green!15, draw=green!70!black, text=green!90!black, font=\bfseries\scriptsize, inner sep=4pt}
  ]
    \node[b1] (ps) at (-4.5, 0) {PHÂN SỐ: $\dfrac{3}{5}$};
    \node[b2] (stp) at (0, 0) {SỐ THẬP PHÂN: $0,6$};
    \node[b3] (pct) at (4.5, 0) {TỈ SỐ \%: $60\%$};

    \draw[teal!85!black, thick, ->] (ps) to[bend left=20] node[above, font=\tiny] {Chia tử cho mẫu} (stp);
    \draw[teal!85!black, thick, ->] (stp) to[bend left=20] node[above, font=\tiny] {$\times 100$ kèm $\%$} (pct);
    \draw[teal!85!black, thick, ->] (pct) to[bend left=20] node[below, font=\tiny] {$60\% = \frac{60}{100}$} (ps);
  \end{tikzpicture}
\end{tcolorbox}

\vspace{3mm}

% --- KHỐI 4: BẢNG GHI NHỚ MỐC QUAN TRỌNG & BẪY TRẮC NGHIỆM ---
\noindent
\begin{minipage}[t]{0.54\textwidth}
  \begin{tcolorbox}[
    enhanced,
    title={\footnotesize\bfseries\color{white}\faBolt\ Các mốc phần trăm vàng cần thuộc lòng},
    colback=green!3!white,
    colframe=green!60!black,
    arc=2.5mm, boxrule=0.8pt, top=1.8mm, bottom=1.8mm, left=2.2mm, right=2.2mm
  ]
    \footnotesize
    $\bullet$ $1 = 100\%$ (Toàn bộ một đơn vị)\\
    $\bullet$ $\dfrac{1}{2} = 0,5 = 50\%$ (Một nửa) \quad $\bullet$ $\dfrac{1}{4} = 0,25 = 25\%$\\
    $\bullet$ $\dfrac{3}{4} = 0,75 = 75\%$ \qquad $\bullet$ $\dfrac{1}{5} = 0,2 = 20\%$\\
    $\bullet$ $\dfrac{1}{10} = 0,1 = 10\%$ \qquad $\bullet$ $\dfrac{1}{8} = 0,125 = 12,5\%$
  \end{tcolorbox}
\end{minipage}%
\hfill
\begin{minipage}[t]{0.44\textwidth}
  \begin{tcolorbox}[
    enhanced,
    title={\footnotesize\bfseries\color{white}\faExclamationTriangle\ Mẹo tránh bẫy trắc nghiệm},
    colback=red!2!white,
    colframe=red!70!black,
    arc=2.5mm, boxrule=0.8pt, top=1.8mm, bottom=1.8mm, left=2.2mm, right=2.2mm
  ]
    \footnotesize
    $\bullet$ Nhầm thứ tự tỉ số: Tỉ số giữa $A$ và $B$ là $\dfrac{A}{B}$, không phải $\dfrac{B}{A}$.\\
    $\bullet$ Khi viết kết quả tỉ số phần trăm: Đã nhân $100$ thì \textbf{bắt buộc phải viết thêm kí hiệu \%} vào bên phải.
  \end{tcolorbox}
\end{minipage}

\end{document}
